\documentclass[12pt]{article}
\usepackage{graphicx}
\usepackage{amsmath}
\usepackage[letterpaper,top=1in,bottom=1in,left=1in,right=1in]{geometry}
\usepackage[utf8]{inputenc}
\usepackage[T1]{fontenc}
\usepackage{textcomp}
\usepackage{array}
\usepackage{tabularx}
\usepackage{booktabs}
\usepackage{calc}
\usepackage{caption}
\usepackage{ragged2e}
\usepackage{etoolbox}
\usepackage{hyperref}
\setlength{\parindent}{0.5in}
\setlength{\parskip}{0pt}
\renewcommand{\arraystretch}{1.06}
\providecommand{\tightlist}{\setlength{\itemsep}{0pt}\setlength{\parskip}{0pt}}
\newcommand{\passthrough}[1]{#1}
\begin{document}

{\centering
{\large\bfseries Supplementary Information}\par\vspace{0.6em}
{\bfseries Ontology-constrained multi-LLM scoring of hypothesis support in the predictive processing literature}\par\vspace{0.6em}
Hamed Nejat, Jacob A. Westerberg, Alexander Maier, Jesse Spencer-Smith, and Andr\'e M. Bastos\par
}
\vspace{0.8em}

\appendix
\section*{Supplementary Information}
\addcontentsline{toc}{section}{Supplementary Information}

\section*{Supplementary Figures}

\subsection*{Supplementary Fig. S1}
\textbf{Supplementary Fig. S1:} Evaluation Pipeline Architecture. The evaluation pipeline flows from left to right in two independent streams (A to B to C and 1 to 2 to 3). The A to B to C pathway establishes the instructions, canonical glossary, and prompt for the models. The 1 to 2 to 3 pathway ingests the paper and combines figures with text to provide a single text object for the models. Steps 5 and 6 loop between models and papers until all evaluations are completed. The final step 7 quantifies and visualizes the evaluations.

\begin{center}
\includegraphics[width=0.86\linewidth,keepaspectratio]{figures/figure2.png}
\end{center}

\subsection*{Supplementary Fig. S2}
\textbf{Supplementary Fig. S2: Pairwise 2D Hypothesis-Space Projections and Human--Council Concordance.} Panels show paper-level hypothesis scores for Local Oddball (LO; left column) and Global Oddball (GO; right column): (A1) LO $H_1 \times H_2$ ($N=20$); (B1) GO $H_1 \times H_2$ ($N=16$); (A2) LO $H_1 \times H_3$ ($N=10$); (B2) GO $H_1 \times H_3$ ($N=7$); (A3) LO $H_2 \times H_3$ ($N=12$); and (B3) GO $H_2 \times H_3$ ($N=10$). Panel $N$ denotes the intersection of papers evaluated by both human experts for both hypotheses shown; per-hypothesis sample sizes are reported in Supplementary Table S14. Blue circles denote Council paper scores and red triangles denote two-expert human-consensus scores. Large gold-bordered markers show panel-specific means $\pm$ SEM across the same independent papers. Dotted gray lines connect paired evaluations for the same study. For comparison, the single-expert evaluation of Westerberg \& Xiong (2025; open marker 30h*) is plotted against its corresponding Council evaluation, demonstrating the selective low-GO outlier displacement.

\begin{center}
\includegraphics[width=0.98\linewidth,keepaspectratio]{figures/figure_S2.pdf}
\end{center}

\subsection*{Supplementary Fig. S3}
\textbf{Supplementary Fig. S3:} Decoding Robustness and Stochastic Replicate Reliability Analysis. (A) Cross-temperature score stability in the primary complete evaluation model (\texttt{gemma-4-31b-it}; 279 calls across 31 papers), showing Pearson correlations ($r \ge 0.9760$) and mean absolute errors ($\text{MAE} \le 0.0421$) across sampling temperatures $T \in \{0.00, 0.35, 0.70\}$. (B) Replicate reliability across independent stochastic generations held at fixed temperatures ($T=0.00: \bar{r}=0.9905, \text{MAE}=0.0126$; $T=0.35: \bar{r}=0.9724, \text{MAE}=0.0384$; $T=0.70: \bar{r}=0.9444, \text{MAE}=0.0585$), confirming that scoring variability across independent stochastic runs remains tightly bounded relative to cross-paper literature variance.

\begin{center}
\includegraphics[width=0.92\linewidth,keepaspectratio]{figures/figure_S_decoding_robustness.png}
\end{center}

\subsection*{Supplementary Fig. S4}
\textbf{Supplementary Fig. S4:} Evidence Availability and Publication-Identity Sensitivity Analysis. Evaluates factor-cell score correspondence between intact manuscript text (FULL) and de-identified text (MASKED, Panel A; Pearson $r = 0.9937, \text{MAE} = 0.0089$ across $N=112$ overlapping factor cells) and between intact manuscript text and restricted text (ABSTRACT Only, Panel B; Pearson $r = 0.7274, \text{MAE} = 0.0758$ across $N=88$ overlapping factor cells), demonstrating that factor scoring is robust to publication-identity removal but sensitive to restriction of available empirical evidence.

\begin{center}
\includegraphics[width=0.92\linewidth,keepaspectratio]{figures/figure_S_evidence_prior.png}
\end{center}

\subsection*{Supplementary Fig. S5}
\textbf{Supplementary Fig. S5:} Figure-Text-Removal Sensitivity Analysis (Full 31-Paper Corpus Ablation). (A) Paper-level mean score correlation across the corpus between intact inputs (FULL) and figure-text-removed inputs (NO\_FIGURE; Spearman $\rho = 0.4718$ across $N=28$ paired empirical paper means; factor-cell level $\text{MAE} = 0.2253$ across all 533 paired cells nested in 31 papers). (B) Factor-cell level directional concordance and stability across the corpus (directional concordance $= 99.42\%$ on non-zero cells [518/521], all-cell sign agreement $= 97.19\%$ [518/533], Pearson $r = 0.8539$), demonstrating that while removing figure descriptions produces measurable paper-level score adjustments, overall factor directionality was preserved across the literature corpus, while paper-level rank correspondence was moderate.

\begin{center}
\includegraphics[width=0.92\linewidth,keepaspectratio]{figures/figure_S_figure_text_removal.png}
\end{center}

\subsection*{Supplementary Fig. S6}
\textbf{Supplementary Fig. S6:} Evaluator Concordance Strata and Single-Study Sensitivity Analysis. (A) Mean squared difference (MSD) across evaluator strata ($K=2$ human domain experts and autonomous LLM council) compared against a null permutation reference distribution ($H_0$, $B=10{,}000$ iterations). Agreement within human experts ($H\text{--}H = 0.1779$) and between human consensus and council ($H\text{--}A = 0.1431$) is significantly tighter than chance ($H_0 = 0.7490, p_{\text{perm}} = 0.0001$, 0/10,000 extreme permutations). (B) Council condition means for Local Oddball (LO) and Global Oddball (GO) across the entire corpus ($N=31$), excluding Westerberg \& Xiong (2025; $N=30$), and for Westerberg \& Xiong alone ($N=1$), demonstrating that the aggregate corpus contrast is robust across published studies and that the preprint represents a selective low-GO outlier.

\begin{center}
\includegraphics[width=0.92\linewidth,keepaspectratio]{figures/figure_S_westerberg_sensitivity.png}
\end{center}

\newpage
\section*{Supplementary Tables}

\begingroup
\small
\noindent\begin{tabularx}{\textwidth}{@{}r p{1.6in} >{\raggedright\arraybackslash}X l@{}}
\toprule
\textbf{\#} & \textbf{Factor Name} & \textbf{Definition} & \textbf{Scope} \\
\midrule
1 & Subtractive Inhibition (SST) & Linear input suppression via somatic inhibition (Somatostatin-mediated). & LO+GO \\
2 & Divisive Inhibition (PV) & Multiplicative gain reduction (Parvalbumin-mediated). & LO+GO \\
3 & Inhibition (GABA) & Chloride-mediated inhibition for prediction suppression. & LO+GO \\
4 & Habituation to Sequence & Habituation suppresses local oddball response relative to a novel local oddball. & LO \\
5 & Synaptic Depression (Adaptation) & Passive fatigue of synaptic efficacy due to repetition of stimulus. & LO \\
6 & Activity Suppression & Reduction in firing rates for expected (predictable) stimuli. & LO+GO \\
7 & Selective Sharpening & Signal-to-noise increase for predictable stimuli, where noise is selectively suppressed to highlight relevant information. & LO+GO \\
8 & Alpha/Beta Mediated Suppression & Low-frequency bands (8--30Hz) associated with predictions, inhibiting prediction error signals. & LO+GO \\
9 & VIP-Mediated Disinhibition & VIP-mediated inhibition of SST/PV neurons, leading to disinhibited pyramidal activity during prediction error. & LO+GO \\
10 & Precision Weighting (Gain) & Top-down amplification of error units, leading to disinhibited gain in response to attended or highly precise stimuli. & LO+GO \\
11 & E/I Balance Shift & Dynamic adjustment toward higher inhibition for predictable stimuli, leading to suppressed neural firing. & LO+GO \\
12 & Omission Response & A generated internal signal due to absence of expected input. & LO+GO \\
\bottomrule
\end{tabularx}
\endgroup
\vspace{0.4em}

\noindent\textbf{Supplementary Table S1:} The factors of Hypothesis 1 (Predictive Suppression). LO = local oddball. GO = global oddball.

\vspace{1.2em}

\begingroup
\small
\noindent\begin{tabularx}{\textwidth}{@{}r p{1.6in} >{\raggedright\arraybackslash}X l@{}}
\toprule
\textbf{\#} & \textbf{Factor Name} & \textbf{Definition} & \textbf{Scope} \\
\midrule
13 & Error Unit Depolarization & Supragranular pyramidal excitation driving ascending error signals. & LO+GO \\
14 & Gamma Oscillations (30--90Hz) & High-frequency rhythms signaling ascending prediction errors. & LO+GO \\
15 & NMDA Receptor Activation & Nonlinear amplification of error signals in superficial cortical layers. & LO+GO \\
16 & Feedforward Laminar Flow & Supragranular-to-granular/supragranular ascending anatomical projection of errors. & LO+GO \\
17 & Top-Down Error Gating & Feedback modulation of ascending prediction error channels. & LO+GO \\
18 & MMN-like Latencies & Early sensory mismatch latencies (100--200ms) characteristic of local error computation. & LO \\
19 & P300-like Latencies & Late cognitive mismatch latencies (300--500ms) characteristic of global contextual error computation. & GO \\
20 & Inter-Laminar Error Divergence & Dissociation of error processing between granular and infragranular layers. & LO+GO \\
21 & Phase-Amplitude Coupling & Theta/alpha phase modulation of gamma error bursts during surprise. & LO+GO \\
22 & Thalamic Error Routing & First-order vs higher-order thalamocortical routing of sensory prediction mismatches. & LO+GO \\
23 & Predictive Error Attenuation & Dampening of ascending prediction errors following successful model updating. & LO+GO \\
24 & State-Dependent Error Routing & Behavioral state and arousal modulation of feedforward error gain. & LO+GO \\
\bottomrule
\end{tabularx}
\endgroup
\vspace{0.4em}

\noindent\textbf{Supplementary Table S2:} The factors of Hypothesis 2 (Feedforward Prediction Errors).

\newpage

\begingroup
\small
\noindent\begin{tabularx}{\textwidth}{@{}r p{1.9in} >{\raggedright\arraybackslash}X l@{}}
\toprule
\textbf{\#} & \textbf{Factor Name} & \textbf{Definition} & \textbf{Scope} \\
\midrule
25 & Cross-Modal Ubiquity & Implementation of predictive coding mechanisms across multiple sensory modalities (visual, auditory, somatosensory). & LO+GO \\
26 & Laminar Invariance of Principles & Preservation of predictive coding laminar logic across sensory, motor, and association cortices. & LO+GO \\
27 & Cross-Species Homology & Conservation of predictive mechanisms across rodent, non-human primate, and human neocortex. & LO+GO \\
28 & Scale-Free Predictive Dynamics & Fractals and hierarchical self-similarity of predictive computations across temporal scales. & LO+GO \\
29 & Universal Error Coding & Common computational principles for error generation irrespective of receptive field complexity. & LO+GO \\
30 & Canonical Microcircuit Generalization & Universal deployment of the canonical cortical microcircuit for inference across areas. & LO+GO \\
31 & Cortico-Subcortical Homology & Shared predictive coding architecture between neocortex and basal ganglia / cerebellum. & LO+GO \\
32 & Neuromodulatory Invariance & Consistent cholinergic/noradrenergic precision tuning across hierarchical stages. & LO+GO \\
33 & Temporal Stability of Ubiquity & The consistent and non-transient presence of these mechanisms across the hierarchy over long recording sessions. & LO+GO \\
34 & Hierarchical Order Stability & Evidence that the predictive mechanism is ubiquitous and not localized to a single hierarchical pole. & LO+GO \\
35 & Population-Wide Ubiquity & Consistency of the mechanism across diverse neural populations and cell types throughout the hierarchy. & LO+GO \\
36 & State-Independent Ubiquity & Presence of the mechanism across different brain states (e.g., wakefulness, sleep, anesthesia) throughout the hierarchy. & LO+GO \\
\bottomrule
\end{tabularx}
\endgroup
\vspace{0.4em}

\noindent\textbf{Supplementary Table S3:} The factors of Hypothesis 3 (Ubiquity).

\vspace{1.5em}

\begingroup
\small
\noindent\begin{tabularx}{\textwidth}{@{}l p{2.8in} >{\raggedright\arraybackslash}X@{}}
\toprule
\textbf{Term / Context} & \textbf{Description} & \textbf{Definition (notation)} \\
\midrule
Oddball & A surprising event, often a mismatch & $E = \{B\}$ given $S = \{A\}$ \\
Local Oddball (LO) & An oddball that is immediate or short-term, often confounded by low-level adaptation mechanisms. & $E = \{AB\}$ given $S = \{AA\}$ \\
Global Oddball (GO) & An oddball that is not LO, often contextual, that is not caused by adaptation mechanisms. & $E = \{AA\}$ given $S = \{AB\}$ \\
\bottomrule
\end{tabularx}
\endgroup
\vspace{0.4em}

\noindent\textbf{Supplementary Table S4:} Context definitions and formal expectation notations.

\newpage

\begingroup
\scriptsize
\renewcommand{\arraystretch}{0.94}
\setlength{\tabcolsep}{2.5pt}
\noindent\begin{tabularx}{\textwidth}{@{}r l p{1.6in} p{1.6in} >{\raggedright\arraybackslash}X@{}}
\toprule
\textbf{\#} & \textbf{Study Identifier} & \textbf{Modality / Preparation} & \textbf{Inclusion Basis} & \textbf{Bibliographic Citation} \\
\midrule
1 & Attinger et al., 2017 & Mouse V1 (2-Photon $\text{Ca}^{2+}$) & Comparative Set (Visuomotor Mismatch) & \emph{Cell} 169, 1291--1302 \\
2 & Bakhtiari et al., 2021 & In Silico (Deep Neural Nets) & Comparative Set (Hierarchical Streams) & \emph{Nat. Commun.} 12, 25164 \\
3 & Bastos et al., 2012 & In Silico / Theory (Microcircuit) & Comparative Set (Canonical Microcircuit) & \emph{Neuron} 74, 695--708 \\
4 & Bastos et al., 2020 & Primate V1/V4/PFC (Laminar ECoG) & Comparative Set (Laminar Rhythms) & \emph{Neuron} 108, 124--137 \\
5 & Bekinschtein et al., 2009 & Human (Scalp EEG) & Comparative Set (LO vs GO Originator) & \emph{PNAS} 106, 1672--1677 \\
6 & Chao et al., 2018 & Primate Temporal/Frontal (ECoG) & Comparative Set (Large-Scale Hierarchy) & \emph{Nat. Commun.} 9, 3188 \\
7 & Friston et al., 2010 & Theory (Free Energy Principle) & Canonical Extension (Bayesian Inference) & \emph{Nat. Rev. Neurosci.} 11, 127 \\
8 & Furutachi et al., 2024 & Mouse V1 (2-Photon / Opto) & Comparative Set (VIP/SST Disinhibition) & \emph{Nature} 633, 398--406 \\
9 & Garrett et al., 2020 & Mouse V1 (2-Photon $\text{Ca}^{2+}$) & Comparative Set (Novelty vs Deviance) & \emph{eLife} 9, e50340 \\
10 & Greedy et al., 2022 & In Silico (Predictive Coding SNN) & Comparative Set (Error Backpropagation) & \emph{NeurIPS} 35, 24213--24225 \\
11 & Hertag et al., 2020 & In Silico (Spiking Interneurons) & Comparative Set (SST/PV Dynamics) & \emph{eLife} 9, e57541 \\
12 & Jiang et al., 2024 & In Silico (Energy-Based Model) & Comparative Set (Sequence Prediction) & \emph{PLOS Comput. Biol.} 20, e1011801 \\
13 & Keller et al., 2012 & Mouse V1 (2-Photon $\text{Ca}^{2+}$) & Comparative Set (Sensorimotor Mismatch) & \emph{Neuron} 74, 809--815 \\
14 & Keller et al., 2018 & Synthesis (Cortical Circuits) & Comparative Set (Circuit Review) & \emph{Neuron} 100, 424--435 \\
15 & Kiebel et al., 2008 & In Silico (Hierarchical Dynamics) & Comparative Set (Temporal Hierarchy) & \emph{PLoS Comput. Biol.} 4, e1000209 \\
16 & Lao et al., 2023 & Human (Auditory Scalp EEG) & Comparative Set (Omission Response) & \emph{Sci. Adv.} 9, eabq8657 \\
17 & Lee et al., 2025 & In Silico (Multi-Area Network) & Comparative Set (Primate Oscillations) & \emph{PLOS Comput. Biol.} 21, e1013469 \\
18 & Mikulasch et al., 2023 & In Silico / Theory (Plasticity) & Comparative Set (Dendritic Error Review) & \emph{Trends Neurosci.} 46, 45--59 \\
19 & Nejad et al., 2025 & Primate / In Silico (Laminar Model) & Comparative Set (Layered Routing) & \emph{J. Neurosci.} 45, 1234 \\
20 & Payeur et al., 2021 & In Silico (Dendritic Plasticity) & Comparative Set (Burst Multiplexing) & \emph{Nat. Neurosci.} 24, 1010--1019 \\
21 & Rao et al., 2024 & Synthesis (Sensorimotor Neocortex) & Canonical Extension (25-Year Synthesis) & \emph{Nat. Neurosci.} 27, 1221--1235 \\
22 & Rao \& Ballard, 1999 & In Silico (Hierarchical Model) & Canonical Extension (Foundational Model) & \emph{Nat. Neurosci.} 2, 79--87 \\
23 & Sacramento et al., 2018 & In Silico (Dendritic Network) & Comparative Set (Somatodendritic Model) & \emph{NeurIPS} 31, 8721--8731 \\
24 & Spratling et al., 2008 & In Silico (Biased Competition) & Comparative Set (PC-BC Synthesis) & \emph{Biol. Cybern.} 98, 283--292 \\
25 & Spratling et al., 2010 & In Silico (V1 Microcircuit) & Comparative Set (V1 Microcircuit Model) & \emph{J. Neurosci.} 30, 3531--3543 \\
26 & Srinivasan et al., 1982 & Fly Retina (Intracellular) & Canonical Extension (Redundancy Baseline) & \emph{Proc. R. Soc. Lond. B} 216, 427 \\
27 & VanDerveer et al., 2021 & Human (Clinical EEG Cohort) & Comparative Set (Sensory Gating Deficits) & \emph{Schizophr. Bull.} 47, 1385 \\
28 & Wacongne et al., 2011 & Human (Auditory MEG) & Comparative Set (Hierarchical Rule MMN) & \emph{PLoS Biol.} 9, e1001189 \\
29 & Wacongne et al., 2012 & Human EEG \& Spiking Network & Comparative Set (Auditory Oddball Model) & \emph{J. Neurosci.} 32, 3665--3678 \\
30 & Westerberg \& Xiong, 2025 & Primate V1/V4/IT (Laminar MUA) & Primary Index Study (MaDeLaNe) & \emph{bioRxiv} 10.1101/2024.10.02 \\
31 & Yamins et al., 2014 & In Silico / Primate IT (HCNN) & Comparative Set (Feedforward Baseline) & \emph{PNAS} 111, 8619--8624 \\
\bottomrule
\end{tabularx}
\endgroup
\vspace{0.4em}

\noindent\textbf{Supplementary Table S5:} 31-Paper Literature Corpus Registry with modality, inclusion basis, and provenance.

\newpage

\begingroup
\small
\noindent\begin{tabularx}{\textwidth}{@{}l c >{\raggedright\arraybackslash}X@{}}
\toprule
\textbf{Model Name} & \textbf{Parameters} & \textbf{Architecture \& Provenance} \\
\midrule
gemma-3-27b-it & 27B & Google DeepMind (Dense Transformer) \\
gemma-4-31b-it & 31B & Google DeepMind (Extended Reasoning Architecture) \\
deepseek-r1-distill-llama-70b & 70B & DeepSeek-AI / Meta Llama Distillation \\
gpt-oss-claude-4.5-sonnet & 20B & Community Distillation / Reasoning \\
mistral-nemo-12b-thinking & 12B & Mistral AI \& NVIDIA \\
olmo-3-32b-think & 32B & Allen Institute for AI (Ai2) \\
phi-4-reasoning-plus & 14B & Microsoft Research \\
qwen3-14b-gemini-3-pro & 14B & Alibaba Cloud / Community Fine-tune \\
qwen3-14b-claude-4.5-sonnet & 14B & Alibaba Cloud / Community Fine-tune \\
qwen3.5-40b-claude-4.5-opus & 40B & Alibaba Cloud / High-Capacity MoE \\
\bottomrule
\end{tabularx}
\endgroup
\vspace{0.4em}

\noindent\textbf{Supplementary Table S6:} Open-weight LLM evaluator models evaluated in this work.

\vspace{1.5em}

\begingroup
\small
\setlength{\tabcolsep}{3.5pt}
\begin{center}
\begin{tabular}{@{}l c c c c c r@{}}
\toprule
\textbf{Evaluator Comparison} & \textbf{Level} & \textbf{Paired $N$} & \textbf{Pearson $r$} & \textbf{$p$} & \textbf{MAE} & \textbf{Concordance} \\
\midrule
Human--Human ($H_1$ vs $H_2$) & Factor cell & 299 & 0.390 & $<0.0001$ & 0.308 & 88.96\% \\
Human Consensus vs LLM Council & Study level & 25 & 0.529 & $0.0125$ & --- & --- \\
Human Consensus vs LLM Council & Factor cell & 175 & 0.162 & $0.0322$ & 0.301 & 93.14\% \\
Council Inter-Model Average & Factor cell & 584 & 0.540 & $<0.0001$ & 0.284 & 85.00\% \\
Shuffled Evidence (Null $H_0$) & Factor cell & 191 & -0.087 & $0.2314$ & 0.662 & 56.54\% \\
\bottomrule
\end{tabular}
\end{center}
\endgroup
\vspace{0.4em}

\noindent\textbf{Supplementary Table S7:} Evaluator Strata Agreement and Error Hierarchy Summary ($K=2$ independent human domain experts). The independent study-level comparison ($N=25$ papers with jointly evaluated factor ratings; Pearson $r=0.529, p=0.0125, 95\%\text{ CI } [0.169, 0.764]$) treats the independent paper as the unit of analysis. Factor-cell comparisons are descriptive summaries across scored cells nested within papers. Directional concordance denotes agreement in sign on non-zero factor cells.

\vspace{1.5em}

\begingroup
\small
\setlength{\tabcolsep}{4pt}
\noindent\begin{tabularx}{\textwidth}{@{}l c c >{\raggedright\arraybackslash}X@{}}
\toprule
\textbf{Adjudication Category} & \textbf{Cells ($N$)} & \textbf{Proportion} & \textbf{Primary Disagreement / Failure Mode} \\
\midrule
Directly supported & 6 & 20.0\% (6/30) & Full primary text, equation, or figure corroboration of Council assignment \\
Partially supported & 18 & 60.0\% (18/30) & Relevant primary evidence present; conservative magnitude attenuation or specificity nuance \\
Unsupported attribution & 6 & 20.0\% (6/30) & Subtype gain conflation ($N=3$), laminar error inversion ($N=1$), untested modality penalty ($N=1$), cross-paper attribution ($N=1$) \\
\midrule
\textbf{Total Audited Disagreements} & \textbf{30} & \textbf{100.0\% (30/30)} & \textbf{Diagnostic audit of $|H - C| \ge 0.50$ cells ($N=296$ eligible paired cells)} \\
\bottomrule
\end{tabularx}
\endgroup
\vspace{0.4em}

\noindent\textbf{Supplementary Table S8:} Diagnostic Primary-Source Audit across 30 Numerically Selected Large Human--Council Disagreement Cells ($|H - C| \ge 0.50$ among 296 eligible paired cells; 10.1\% of paired universe). Evaluates whether large disagreements reflect unsupported evidence attribution or differences in magnitude, contextual interpretation, or biological specificity despite relevant primary-source evidence. Because this diagnostic sample was enriched for large discrepancies, these proportions characterize disagreement modes and do not represent corpus-wide accuracy or hallucination rates.

\newpage

\begingroup
\small
\setlength{\tabcolsep}{4pt}
\noindent\begin{tabularx}{\textwidth}{@{}p{2.2in} c c c >{\raggedright\arraybackslash}X@{}}
\toprule
\textbf{Decoding Setting / Comparison} & \textbf{Evaluated Units ($N$)} & \textbf{Correlation} & \textbf{MAE} & \textbf{Paper Ranking Stability} \\
\midrule
\textbf{Temperature Stability (\texttt{gemma-4-31b-it})} & & & & \\
\quad $T = 0.00$ vs. $T = 0.35$ & 679 factor cells & Pearson $r = 0.9900$ & $0.0266$ & Spearman $\rho_{\text{rank}} = 0.9596$ \\
\quad $T = 0.00$ vs. $T = 0.70$ & 679 factor cells & Pearson $r = 0.9760$ & $0.0421$ & Spearman $\rho_{\text{rank}} = 0.9212$ \\
\quad $T = 0.35$ vs. $T = 0.70$ & 679 factor cells & Pearson $r = 0.9816$ & $0.0367$ & Spearman $\rho_{\text{rank}} = 0.9482$ \\
\midrule
\textbf{Replicate Reliability (\texttt{gemma-4-31b-it})} & & & & \\
\quad Replicates at $T = 0.00$ & 639 triples & Mean $\bar{r} = 0.9905$ & $0.0126$ & Spearman $\rho_{\text{rank}} \ge 0.9818$ \\
\quad Replicates at $T = 0.35$ & 619 triples & Mean $\bar{r} = 0.9724$ & $0.0384$ & Spearman $\rho_{\text{rank}} \ge 0.9782$ \\
\quad Replicates at $T = 0.70$ & 572 triples & Mean $\bar{r} = 0.9444$ & $0.0585$ & Spearman $\rho_{\text{rank}} \ge 0.9745$ \\
\bottomrule
\end{tabularx}
\endgroup
\vspace{0.4em}

\noindent\textbf{Supplementary Table S9:} Factorial Robustness Analysis across Sampling Temperatures and Stochastic Replicates for the Complete Evaluation Model (\texttt{gemma-4-31b-it}, 279/279 factorial conditions). Quantitative factorial statistics are reported exclusively for \texttt{gemma-4-31b-it} because it completed full factorial coverage across all 31 papers $\times$ 3 temperatures $\times$ 3 repeats. Extended chain-of-thought models (\texttt{phi-4-reasoning-plus} and \texttt{olmo-3-32b-think}) completed 130/279 and 30/279 calls, respectively, but frequently exhausted dynamic generation token budgets during chain-of-thought reasoning prior to emitting the terminal JSON scoring block, precluding equivalent comparative factorial statistics. Paper ranking stability measures Spearman correlation between paper mean scores across conditions ($N=29$ evaluable papers).

\vspace{1.5em}

\begingroup
\small
\noindent\textbf{Panel A: Paired Score Sensitivity to Identity Removal and Evidence Restriction}
\vspace{0.2em}

\noindent\begin{tabularx}{\textwidth}{@{}l c c c >{\raggedright\arraybackslash}X@{}}
\toprule
\textbf{Comparison} & \textbf{Paired Cells ($N$)} & \textbf{Pearson $r$} & \textbf{MAE} & \textbf{Scientific Interpretation} \\
\midrule
FULL vs MASKED Text & 112 & 0.9937 & 0.0089 & Tests dependence on explicit publication identity \\
FULL vs ABSTRACT Only & 88 & 0.7274 & 0.0758 & Tests sensitivity to restriction of available paper evidence \\
\bottomrule
\end{tabularx}

\vspace{0.8em}
\noindent\textbf{Panel B: Abstract-Only Sensitivity by Hypothesis Family}
\vspace{0.2em}

\noindent\begin{tabularx}{\textwidth}{@{}l c c c >{\raggedright\arraybackslash}X@{}}
\toprule
\textbf{Hypothesis Family} & \textbf{Paired Cells ($N$)} & \textbf{Pearson $r$} & \textbf{MAE} & \textbf{Coverage Note} \\
\midrule
Hypothesis 1 (Predictive Suppression) & 30 & 0.6003 & 0.1444 & 5 papers with evaluable abstract evidence \\
Hypothesis 2 (Feedforward Error) & 38 & 0.6776 & 0.0614 & 4 papers with evaluable abstract evidence \\
Hypothesis 3 (Ubiquity) & 20 & 1.0000 & 0.0000 & 2 papers; insufficient coverage for general inference \\
\bottomrule
\end{tabularx}

\vspace{0.8em}
\noindent\textbf{Panel C: Zero-Manuscript-Evidence Control (Prior Only)}
\vspace{0.2em}

\noindent\begin{tabularx}{\textwidth}{@{}l c c c >{\raggedright\arraybackslash}X@{}}
\toprule
\textbf{Input Condition} & \textbf{Calls ($N$)} & \textbf{LO Mean [$\pm$SD]} & \textbf{GO Mean [$\pm$SD]} & \textbf{Context Contrast ($\Delta_{\text{LO}-\text{GO}}$)} \\
\midrule
Intact Evidence (FULL Reference) & 18 & $0.512 \pm 0.081$ & $0.474 \pm 0.076$ & $+0.038$ (Active context differentiation) \\
Prior Only (Zero Manuscript Evidence) & 18 & $0.450 \pm 0.110$ & $0.450 \pm 0.108$ & $\phantom{+}0.000$ (Parity across contexts) \\
\bottomrule
\end{tabularx}
\endgroup
\vspace{0.4em}

\noindent\textbf{Supplementary Table S10:} Evidence Availability, Publication-Identity, and Prior-Only Sensitivity Analyses. De-identifying full manuscript text produced minimal score adjustments (FULL vs MASKED, Panel A), whereas restricting available evidence to titles and abstracts produced larger deviations (FULL vs ABSTRACT, Panel A) with the largest attenuation in predictive suppression mechanisms (Panel B). In the zero-manuscript-evidence condition (Panel C), models retained baseline affirmative support from the ontology and task framing, but the Local vs. Global Oddball context differentiation observed under intact manuscript evidence was completely absent. Abstract-only Hypothesis 3 coverage was limited to 20 paired cells from two papers and is therefore not interpreted as evidence of general invariance across the literature.

\vspace{1.5em}

\begingroup
\small
\setlength{\tabcolsep}{4pt}
\noindent\begin{tabularx}{\textwidth}{@{}l c c c >{\raggedright\arraybackslash}X@{}}
\toprule
\textbf{Ablation Condition} & \textbf{Paired Cells ($N$)} & \textbf{Pearson $r$} & \textbf{MAE} & \textbf{Directional Concordance} \\
\midrule
FULL Baseline (With Figures) & 533 & 1.000 & 0.000 & 100.0\% (Baseline) \\
NO\_FIGURE (Figures Stripped) & 533 & 0.8539 & 0.2253 & 99.42\% (518/521 non-zero cells) \\
\bottomrule
\end{tabularx}
\endgroup
\vspace{0.4em}

\noindent\textbf{Supplementary Table S11:} Figure-Text-Removal Sensitivity Summary ($N=31$ papers).

\newpage

\begingroup
\small
\setlength{\tabcolsep}{4pt}
\noindent\begin{tabularx}{\textwidth}{@{}l >{\raggedright\arraybackslash}X c c c@{}}
\toprule
\textbf{Dimension} & \textbf{Stratum} & \textbf{Papers ($N$)} & \textbf{Council LO Mean} & \textbf{Council GO Mean} \\
\midrule
Study Type & Empirical & 13 & $0.419$ & $0.428$ \\
Study Type & Computational / Theory & 18 & $0.351$ & $0.365$ \\
Modality & In Silico & 19 & $0.363$ & $0.376$ \\
Modality & Human In Vivo & 7 & $0.357$ & $0.416$ \\
Modality & Animal In Vivo & 5 & $0.511$ & $0.424$ \\
Publication Era & 1982--2014 & 12 & $0.368$ & $0.393$ \\
Publication Era & 2015--2020 & 7 & $0.447$ & $0.488$ \\
Publication Era & 2021--2025 & 12 & $0.347$ & $0.309$ \\
Publication Status & Peer-Reviewed Journal & 30 & $0.381$ & $0.392$ \\
Publication Status & Preprint (Westerberg \& Xiong 2025) & 1 & $0.166$ & $-0.244$ \\
\bottomrule
\end{tabularx}
\endgroup
\vspace{0.4em}

\noindent\textbf{Supplementary Table S12:} Corpus Subgroup Sensitivity Analysis across Taxonomic Strata. Demonstrates that council-estimated evidence support across Local and Global Oddball contexts varies by experimental modality and publication era, with Westerberg \& Xiong (2025) exhibiting a distinct negative Global Oddball profile.

\vspace{1.5em}

\begingroup
\small
\setlength{\tabcolsep}{3pt}
\noindent\begin{tabularx}{\textwidth}{@{}p{1.7in} c c c >{\centering\arraybackslash}X@{}}
\toprule
\textbf{Validation Axis} & \textbf{Evaluated ($N$)} & \textbf{Strictly Correct} & \textbf{Partial / Minor} & \textbf{Audit Pass Rate} \\
\midrule
Axis A: Numerical Trend \& Data Extraction & 10 & 9 & 1 (Minor phrasing) & 100.0\% PASS ($10/10$) \\
Axis B: Empirical Effect Direction & 10 & 10 & 0 & 100.0\% PASS ($10/10$) \\
Axis C: HPC-36 Ontology Factor Mapping & 10 & 9 & 1 (Invertebrate) & 100.0\% PASS ($10/10$) \\
\midrule
\textbf{Overall Manual Validation Outcome} & \textbf{10 Figures} & \textbf{10/10} & \textbf{0 Fatal Inversions} & \textbf{100.0\% PASS} \\
\bottomrule
\end{tabularx}
\endgroup
\vspace{0.4em}

\noindent\textbf{Supplementary Table S13:} Manual Expert Validation Summary of Generated Multimodal Figure Descriptions ($N=10$ representative figures audited against primary source PDFs). All 10 figures met the prespecified audit criterion (Audit Pass Rate: 100.0\%); minor phrasing and partial ontology mappings are reported separately.

\vspace{1.5em}

\begingroup
\footnotesize
\setlength{\tabcolsep}{2.5pt}
\noindent\begin{tabularx}{\textwidth}{@{}l c c c c >{\centering\arraybackslash}X r@{}}
\toprule
\textbf{Context} & \textbf{Hyp.} & \textbf{$N$} & \textbf{Human Mean} & \textbf{Council Mean} & \textbf{Human$-$Council [95\% CI]} & \textbf{$t(df), p$} \\
\midrule
\textbf{Local Oddball}  & $H_1$ & 21 & $+0.537 \pm 0.041$ & $+0.337 \pm 0.027$ & $+0.200$ $[0.111, 0.289]$ & $t(20)=4.69, p=0.0001$ \\
\textbf{(LO)}           & $H_2$ & 23 & $+0.552 \pm 0.045$ & $+0.363 \pm 0.026$ & $+0.189$ $[0.117, 0.261]$ & $t(22)=5.42, p<0.0001$ \\
                        & $H_3$ & 14 & $+0.520 \pm 0.034$ & $+0.268 \pm 0.045$ & $+0.252$ $[0.161, 0.344]$ & $t(13)=5.95, p<0.0001$ \\
\midrule
\textbf{Global Oddball} & $H_1$ & 18 & $+0.600 \pm 0.048$ & $+0.299 \pm 0.031$ & $+0.302$ $[0.182, 0.422]$ & $t(17)=5.32, p<0.0001$ \\
\textbf{(GO)}           & $H_2$ & 22 & $+0.560 \pm 0.043$ & $+0.259 \pm 0.049$ & $+0.300$ $[0.205, 0.396]$ & $t(21)=6.54, p<0.0001$ \\
                        & $H_3$ & 11 & $+0.557 \pm 0.065$ & $+0.312 \pm 0.057$ & $+0.244$ $[0.080, 0.408]$ & $t(10)=3.32, p=0.0077$ \\
\bottomrule
\end{tabularx}
\endgroup
\vspace{0.4em}

\noindent\textbf{Supplementary Table S14:} Primary Paper-Level Hypothesis Calibration and Concordance ($K=2$ Human Consensus vs. Autonomous Council). Analysis is restricted to independent papers where both human experts ($K=2$) evaluated mechanisms within the specific hypothesis. Council models were averaged within paper first. SEM denotes standard error across independent papers ($\text{SD}/\sqrt{N}$). Paired differences (Human$-$Council mean difference) evaluate systematic score attenuation. Paired $t$-tests ($t(df)$) and two-sided $p$-values evaluate the difference in mean support.

\vspace{1.5em}

\begingroup
\small
\setlength{\tabcolsep}{4.5pt}
\noindent\begin{tabularx}{\textwidth}{@{}p{2.1in} c c c c c r@{}}
\toprule
\textbf{Omitted Evaluator Model} & \textbf{Paired $N$} & \textbf{Pearson $r$} & \textbf{$p$} & \textbf{Spearman $\rho$} & \textbf{$p$} & \textbf{MAE} \\
\midrule
\texttt{deepseek-r1-distill-llama-70b} & 1710 & 0.9553 & $<0.0001$ & 0.9589 & $<0.0001$ & 0.0465 \\
\texttt{gemma-3-27b}                   & 1733 & 0.9747 & $<0.0001$ & 0.9472 & $<0.0001$ & 0.0408 \\
\texttt{gpt-oss-claude-4.5-sonnet}     & 1801 & 0.9920 & $<0.0001$ & 0.9844 & $<0.0001$ & 0.0086 \\
\texttt{mistral-nemo-12b-thinking}     & 1775 & 0.9732 & $<0.0001$ & 0.9698 & $<0.0001$ & 0.0239 \\
\texttt{olmo-3-32b-think}              & 1798 & 0.9900 & $<0.0001$ & 0.9800 & $<0.0001$ & 0.0122 \\
\texttt{phi-4-reasoning-plus}          & 1778 & 0.9825 & $<0.0001$ & 0.9735 & $<0.0001$ & 0.0255 \\
\texttt{qwen3-14b-gemini-3-pro}        & 1800 & 0.9797 & $<0.0001$ & 0.9649 & $<0.0001$ & 0.0243 \\
\texttt{qwen3.5-40b-claude-4.5-opus}   & 1724 & 0.9497 & $<0.0001$ & 0.9178 & $<0.0001$ & 0.0440 \\
\midrule
\multicolumn{7}{@{}l}{\textit{Evaluator profiles with hypothesis-level scores but unassigned individual factor cells:}} \\
\texttt{gemma-4-31b}                   & 1816 & --- & --- & --- & --- & --- \\
\texttt{gpt-oss-safeguard-120b}        & 1816 & --- & --- & --- & --- & --- \\
\bottomrule
\end{tabularx}
\endgroup
\vspace{0.4em}

\noindent\textbf{Supplementary Table S15:} Leave-One-Model-Out Council Score Stability Analysis across Evaluator Architectures. Evaluates the sensitivity of aggregate Council scores to single-model removal by comparing the complete Council continuous mean against the Council recomputed after omission of each contributing evaluator across matched factor cells ($N = 1710\text{--}1801$). Across all active evaluator omissions ($N=8$ model omissions), Council scores remained highly stable: mean Pearson $r = 0.9747$ (95\% CI: $[0.9619, 0.9874]$, range $0.9497\text{--}0.9920$), mean Spearman $\rho = 0.9621$ (95\% CI: $[0.9441, 0.9800]$, range $0.9178\text{--}0.9844$), and mean $\text{MAE} = 0.0282$ (95\% CI: $[0.0163, 0.0402]$, range $0.0086\text{--}0.0465$). Confidence intervals are computed using Student's $t$-distribution ($df = N_{\text{omissions}}-1 = 7$) with model omission as the unit. MAE is reported in original $[-1.0, +1.0]$ score units. Two evaluator profiles contributed hypothesis-level Council scores but no individual factor-cell scores and therefore were not evaluable in this factor-level leave-one-out analysis.

\end{document}
